\documentclass[11pt,a4paper]{article}
\usepackage[utf8]{inputenc}
\usepackage[T1]{fontenc}
\usepackage[brazil]{babel}
\usepackage{lmodern,microtype}
\usepackage{geometry}
\geometry{top=2.4cm,bottom=2.4cm,left=2.8cm,right=2.4cm}
\usepackage{setspace}
\setstretch{1.18}
\usepackage{amsmath,amssymb}
\usepackage{booktabs,tabularx,array}
\usepackage{siunitx}
\sisetup{output-decimal-marker={,}}
\usepackage{xcolor}
\definecolor{azul}{HTML}{164E8A}
\definecolor{laranja}{HTML}{B45309}
\definecolor{verde}{HTML}{047857}
\definecolor{vermelho}{HTML}{B91C1C}
\usepackage{tikz}
\usetikzlibrary{arrows.meta,positioning,fit}
\usepackage{pgfplots}
\usepgfplotslibrary{groupplots}
\pgfplotsset{compat=1.18}
\usepackage{caption}
\captionsetup{font=small,labelfont=bf}
\usepackage{enumitem}
\setlist{nosep,leftmargin=1.5em}
\usepackage[hidelinks]{hyperref}
\hypersetup{pdftitle={Equipe 04 - Modulacao em amplitude - Capitulo 5},pdfauthor={Gabriela Pontes; Julia Goncalves; Placido Cordeiro}}
\newif\ifcapturavhil
\IfFileExists{resultados_vhil/medidas.tex}{\capturavhiltrue}{\capturavhilfalse}
\newcommand{\arquivo}[1]{\texttt{\detokenize{#1}}}
\newcommand{\statuscaptura}{\ifcapturavhil Resultados experimentais VHIL incluídos\else Prévia técnica: captura VHIL ainda não incorporada\fi}

\begin{document}
\begin{titlepage}
\thispagestyle{empty}
\begin{center}
{\large UNIVERSIDADE FEDERAL DE ALAGOAS\par}
\vspace{0.15cm}
{\large INSTRUMENTAÇÃO ELETRÔNICA -- ECOM060\par}
\vfill
{\LARGE\bfseries MODULAÇÃO EM AMPLITUDE\par}
\vspace{0.2cm}
{\Large Equipe 04 -- Capítulo 5, seção 5.2\par}
\vspace{0.35cm}
{\large Teoria, implementação no TyphoonSim e análise espectral\par}
\vspace{0.8cm}
{\normalsize\color{\ifcapturavhil verde\else vermelho\fi}\bfseries\statuscaptura\par}
\vfill
Gabriela Pontes\\
Júlia Gonçalves\\
Plácido Cordeiro
\vfill
\begin{minipage}{0.76\textwidth}
\small
Relatório da atividade distribuída do Capítulo 5 de \emph{Fundamentos de Instrumentação}, de Luís Antonio Aguirre. O estudo implementa modulação AM em TyphoonSim/Virtual HIL e compara profundidades de modulação abaixo, no limite e acima da unidade.
\par\medskip
Professor: Maurício Beltrão Rossiter
\end{minipage}
\vfill
Maceió -- AL\\2026
\end{center}
\end{titlepage}

\pagenumbering{roman}
\section*{Resumo}
Este relatório apresenta a modulação em amplitude (AM) de portadora presente para um tom modulante. Partindo de uma mensagem senoidal normalizada de \SI{25}{\hertz} e de uma portadora de \SI{500}{\hertz}, demonstra-se que o produto $s(t)=[1+\mu m(t)]c(t)$ produz uma linha espectral central e duas bandas laterais em \SI{475}{\hertz} e \SI{525}{\hertz}. Foram definidos três índices, $\mu=0{,}5$, $1{,}0$ e $1{,}5$, para examinar submodulação, modulação crítica e sobremodulação. As curvas analíticas permitem antecipar a amplitude das linhas e a inversão de fase quando a envoltória assinada fica negativa. O modelo TyphoonSim e a aquisição automática de cinco probes foram preparados com passo de \SI{100}{\micro\second}.
\ifcapturavhil
 A seção experimental usa a captura VHIL anexada, com tabela, formas de onda, FFT e hashes de rastreabilidade.
\else
 Esta versão contém somente previsões analíticas: não havia Typhoon HIL instalado no computador de autoria, portanto a captura VHIL precisa ser executada antes da entrega ao professor. Nenhuma curva teórica é apresentada como medição.
\fi

\noindent\textbf{Palavras-chave:} modulação em amplitude; bandas laterais; índice de modulação; TyphoonSim; análise espectral.

\clearpage
\tableofcontents
\clearpage
\pagenumbering{arabic}

\section{Introdução e objetivos}
Na modulação em amplitude, uma grandeza de interesse controla a amplitude de uma portadora mais rápida. Isso desloca informação de baixa frequência para uma vizinhança espectral da frequência da portadora, situação importante em instrumentação e em sistemas de comunicação. A Equipe 04 recebeu a seção 5.2 do Capítulo 5 de Aguirre, com a incumbência de explicar a teoria e demonstrá-la no TyphoonSim. A orientação exige forma de onda temporal, bandas laterais e comparação entre índices de modulação. Como não haverá exposição oral, apresentam-se as definições, a derivação, as escolhas numéricas, o procedimento de captura e os critérios de interpretação.

O objetivo é verificar a correspondência entre (i) o sinal AM produzido pelo modelo, (ii) as três frequências previstas pela expansão trigonométrica e (iii) o comportamento da envoltória nos regimes $\mu<1$, $\mu=1$ e $\mu>1$. O modelo não contém demodulador: a demodulação constitui o subtema atribuído à Equipe 05.

\section{Fundamentação teórica}
\subsection{Equação da AM de portadora presente}
Adota-se uma mensagem normalizada $m(t)=\cos(2\pi f_m t)$, com $|m(t)|\leq 1$, e uma portadora $c(t)=A_c\cos(2\pi f_c t)$. A modulação em amplitude convencional é
\begin{equation}
 u(t)=1+\mu m(t),\qquad s(t)=u(t)c(t)
 =A_c[1+\mu\cos(2\pi f_m t)]\cos(2\pi f_c t),
 \label{eq:am}
\end{equation}
em que $A_c$ é a amplitude de pico da portadora e $\mu\geq0$ é o índice de modulação. $u(t)$ é a \emph{envoltória assinada}; a envoltória geométrica superior do sinal observado é $A_c|u(t)|$. Essa distinção é essencial quando $u(t)$ cruza zero.

Aplicando $\cos a\cos b=\tfrac12[\cos(a+b)+\cos(a-b)]$ à Equação~\eqref{eq:am}, obtém-se
\begin{equation}
 s(t)=A_c\cos(2\pi f_c t)
 +\frac{A_c\mu}{2}\cos[2\pi(f_c-f_m)t]
 +\frac{A_c\mu}{2}\cos[2\pi(f_c+f_m)t].
 \label{eq:spectrum}
\end{equation}
Para um único tom, aparecem linhas em $f_c-f_m$, $f_c$ e $f_c+f_m$. Cada banda lateral tem amplitude de pico $A_c\mu/2$; a largura ocupada entre as bandas laterais é $2f_m$. O espectro bilateral possui também as linhas correspondentes em frequências negativas, omitidas nos gráficos unilaterais.

\subsection{Profundidade de modulação e envoltória}
Como $\min m(t)=-1$, o valor mínimo da envoltória assinada é $u_{\min}=1-\mu$. A Tabela~\ref{tab:regimes} resume os três casos. Para $\mu>1$, $u(t)<0$ em parte do ciclo: como $u(t)c(t)=-|u(t)|c(t)$ nesse intervalo, a portadora sofre inversão de fase de $\pi$ rad. Um detector de envoltória simples observa $|u(t)|$, não $u(t)$, e por isso deixa de reproduzir fielmente a mensagem nesse regime.

Quando $0\leq\mu\leq1$, os extremos da envoltória geométrica são $E_{\max}=A_c(1+\mu)$ e $E_{\min}=A_c(1-\mu)$. Logo, $\mu=(E_{\max}-E_{\min})/(E_{\max}+E_{\min})$. Esta fórmula de extremos deixa de distinguir corretamente a sobremodulação: após o dobramento por valor absoluto, o mínimo geométrico passa a zero, embora a envoltória \emph{assinada} continue abaixo de zero. Por isso o relatório também inspeciona o probe $u(t)$ e estima $\mu$ pela razão entre bandas laterais.

\begin{table}[htbp]
\centering
\caption{Critério físico para classificar a profundidade de modulação.}
\label{tab:regimes}
\begin{tabularx}{\textwidth}{@{}l c X@{}}
\toprule
Regime & Índice & Consequência para $u(t)=1+\mu m(t)$\\
\midrule
Submodulação & $0<\mu<1$ & $u_{\min}>0$; a envoltória não toca zero.\\
Limite de 100\% & $\mu=1$ & $u_{\min}=0$; a envoltória toca zero.\\
Sobremodulação & $\mu>1$ & $u_{\min}<0$; há cruzamentos por zero e inversões de fase.\\
\bottomrule
\end{tabularx}
\end{table}

\textbf{Nota importante:} no multiplicador ideal da Equação~\eqref{eq:am}, aumentar $\mu$ acima de 1 não cria, por si só, novas frequências no sinal AM. As duas bandas laterais crescem linearmente. A distorção não linear surge se se tenta recuperar a mensagem tomando o módulo da envoltória ou por um detector de envoltória convencional.

\section{Modelo implementado no TyphoonSim}
\begin{figure}[htbp]
\centering
\begin{tikzpicture}[node distance=1.1cm and 1.0cm,>=Latex,
box/.style={draw=azul,rounded corners=2pt,fill=blue!4,align=center,minimum height=1.25cm,text width=3.15cm},
smallbox/.style={draw=laranja,rounded corners=2pt,fill=orange!5,align=center,minimum height=1.0cm,text width=2.5cm}]
\node[box] (src) {Fontes senoidais\\$m(t)$, $c(t)$};
\node[smallbox,below=of src] (seq) {Sequenciador\\$\mu=0{,}5;\ 1{,}0;\ 1{,}5$};
\node[box,right=1.35cm of src,yshift=-0.6cm] (mod) {Modulador AM\\$u=1+\mu m$\\$s=uc$};
\node[box,right=1.35cm of mod] (probes) {Cinco probes\\mensagem, portadora,\\índice, $u$ e $s$};
\draw[->,thick] (src.east) -- node[above,sloped,scale=.78] {$m,c$} (mod.west);
\draw[->,thick] (seq.east) -| node[below,near start,scale=.78] {$\mu$} (mod.south);
\draw[->,thick] (mod.east) -- (probes.west);
\end{tikzpicture}
\caption{Arquitetura do arquivo \arquivo{Equipe_04_ModulacaoAM.tse}. As saídas alimentam o Capture VHIL; não se utiliza leitura pontual para construir o espectro.}
\label{fig:blocos}
\end{figure}

O esquemático usa um bloco de fontes/sequenciamento e um bloco modulador, ambos \emph{C function} genéricos, conectados a probes públicos. Não usa chaves ou diodos de potência, evitando a incompatibilidade com HIL404 Base 1 encontrada em uma atividade anterior. O tipo C \texttt{double} é empregado nas variáveis e constantes internas; a palavra \texttt{real} aparece apenas na declaração dos terminais Typhoon. Não há regulador, conversor Buck nem malha de realimentação: esta é uma atividade nova, independente daquela planta.

\begin{table}[htbp]
\centering
\caption{Parâmetros declarados no modelo e previstos para a captura.}
\label{tab:parametros}
\begin{tabular}{@{}l l l@{}}
\toprule
Grandeza & Valor & Justificativa\\
\midrule
Frequência da mensagem $f_m$ & \SI{25}{\hertz} & Período de \SI{40}{\milli\second}.\\
Frequência da portadora $f_c$ & \SI{500}{\hertz} & 20 amostras por ciclo.\\
Amplitude da portadora $A_c$ & \SI{1}{\volt} de pico & Facilita comparação espectral.\\
Passo do sinal $T_s$ & \SI{100}{\micro\second} & Taxa $f_s=\SI{10}{\kilo\hertz}$.\\
Índices $\mu$ & $0{,}5;\ 1{,}0;\ 1{,}5$ & Regimes abaixo, no e acima de 100\%.\\
Duração de cada regime & \SI{0.4}{\second} & 10 ciclos de $m$, 200 de $c$.\\
Captura solicitada & 36000 pontos & \SI{3.6}{\second}, três ciclos de regimes.\\
\bottomrule
\end{tabular}
\end{table}

O sequenciador percorre 4000 passos em cada índice e repete a sequência. Assim, mesmo que o Capture inicie no meio de um regime, uma janela de 36000 pontos contém uma passagem completa por todos eles. Com $f_s=\SI{10}{\kilo\hertz}$, a frequência de Nyquist é \SI{5}{\kilo\hertz}, muito acima da banda lateral superior de \SI{525}{\hertz}. O ensaio é inteiramente de processamento de sinais; não representa um transmissor de RF físico.

\section{Procedimento de aquisição e análise}
No computador que contém Typhoon HIL Control Center, o script \arquivo{capturar_vhil.py} abre e compila o \arquivo{.tse}, carrega o \arquivo{.cpd} em TyphoonSim/Virtual HIL e inicia uma captura forçada dos cinco probes. A chamada \texttt{hil.start\_capture} solicita 36000 amostras. A decimação inteira é calculada a partir do passo-base $T_{\mathrm{base}}$ informado pelo modelo carregado, de modo que $D\,T_{\mathrm{base}}\simeq\SI{100}{\micro\second}$; por exemplo, $T_{\mathrm{base}}=\SI{0.2}{\micro\second}$ requer $D=500$. O vetor de tempo e a matriz de canais vêm da API do Typhoon. O script confere taxa efetiva, uniformidade temporal, presença dos três índices e as identidades $u=1+\mu m$ e $s=uc$ antes de gravar CSV e metadados. Os hashes SHA-256 do modelo, do compilado e do CSV preservam a rastreabilidade.

O script \arquivo{analisar_captura.py} seleciona uma janela contígua de cada regime e calcula duas representações complementares. Para os gráficos espectrais, aplica uma janela de Hann e transforma a série amostrada para obter amplitudes unilaterais; a janela reduz vazamento mas alarga os picos. Para a tabela quantitativa, ajusta simultaneamente seno e cosseno nas frequências \SI{475}{\hertz}, \SI{500}{\hertz} e \SI{525}{\hertz} por mínimos quadrados. O módulo de cada par de coeficientes fornece a amplitude de pico. O índice recuperado do espectro é
\begin{equation}
 \widehat\mu=\frac{\widehat A_{475}+\widehat A_{525}}{\widehat A_{500}}.
 \label{eq:mu_est}
\end{equation}
O ajuste também fornece o RMS do resíduo. A função de análise recusa dados incompletos ou sem metadados VHIL e não mistura curvas teóricas com arquivos medidos.

\section{Resultados previstos analiticamente}
Com os valores da Tabela~\ref{tab:parametros}, a Equação~\eqref{eq:spectrum} prevê banda lateral inferior em \SI{475}{\hertz}, portadora em \SI{500}{\hertz} e banda lateral superior em \SI{525}{\hertz}. A largura entre as bandas laterais é \SI{50}{\hertz}. As amplitudes esperadas estão na Tabela~\ref{tab:teoria}. Elas são \emph{previsões calculadas}, não saídas do TyphoonSim.

\begin{table}[htbp]
\centering
\caption{Previsão ideal para os três índices. Amplitudes unilaterais de pico, $A_c=\SI{1}{\volt}$.}
\label{tab:teoria}
\begin{tabular}{@{}lrrrr@{}}
\toprule
$\mu$ & $A_{475}$ [V] & $A_{500}$ [V] & $A_{525}$ [V] & $u_{\min}$ [V]\\
\midrule
$0{,}5$ & $0{,}25$ & $1{,}00$ & $0{,}25$ & $0{,}50$\\
$1{,}0$ & $0{,}50$ & $1{,}00$ & $0{,}50$ & $0{,}00$\\
$1{,}5$ & $0{,}75$ & $1{,}00$ & $0{,}75$ & $-0{,}50$\\
\bottomrule
\end{tabular}
\end{table}

\begin{figure}[p]
\centering
\begin{tikzpicture}
\begin{groupplot}[
 group style={group size=1 by 3,vertical sep=0.9cm},
 width=.88\textwidth,height=5.0cm,
 xmin=0,xmax=40,ymin=-2.7,ymax=2.7,
 xtick={0,10,20,30,40},grid=both,
 tick label style={font=\scriptsize},label style={font=\small},title style={font=\small},
 ylabel={Tensão (V)},xlabel={Tempo (ms)}]
\nextgroupplot[title={$\mu=0{,}5$: envoltória estritamente positiva}]
\addplot[azul,thin] table[x=tempo_ms,y=sinal_am,col sep=comma]{referencia_analitica/tempo_teorico_mu05.csv};
\addplot[laranja,thick] table[x=tempo_ms,y=env_pos,col sep=comma]{referencia_analitica/tempo_teorico_mu05.csv};
\addplot[laranja,thick] table[x=tempo_ms,y=env_neg,col sep=comma]{referencia_analitica/tempo_teorico_mu05.csv};
\nextgroupplot[title={$\mu=1{,}0$: envoltória toca zero}]
\addplot[azul,thin] table[x=tempo_ms,y=sinal_am,col sep=comma]{referencia_analitica/tempo_teorico_mu10.csv};
\addplot[laranja,thick] table[x=tempo_ms,y=env_pos,col sep=comma]{referencia_analitica/tempo_teorico_mu10.csv};
\addplot[laranja,thick] table[x=tempo_ms,y=env_neg,col sep=comma]{referencia_analitica/tempo_teorico_mu10.csv};
\nextgroupplot[title={$\mu=1{,}5$: envoltória assinada negativa perto de 20 ms}]
\addplot[azul,thin] table[x=tempo_ms,y=sinal_am,col sep=comma]{referencia_analitica/tempo_teorico_mu15.csv};
\addplot[laranja,thick] table[x=tempo_ms,y=env_pos,col sep=comma]{referencia_analitica/tempo_teorico_mu15.csv};
\addplot[laranja,thick] table[x=tempo_ms,y=env_neg,col sep=comma]{referencia_analitica/tempo_teorico_mu15.csv};
\addplot[vermelho,densely dashed,thick] table[x=tempo_ms,y=env_abs,col sep=comma]{referencia_analitica/tempo_teorico_mu15.csv};
\end{groupplot}
\end{tikzpicture}
\caption{Curvas \textbf{teóricas}, calculadas diretamente da Equação~\eqref{eq:am}, para um período da mensagem. Azul: sinal AM; laranja: $\pm u(t)$; vermelho tracejado, apenas no último painel: $|u(t)|$, envoltória geométrica superior que um detector simples seguiria. Não são capturas VHIL.}
\label{fig:tempo_teoria}
\end{figure}

\begin{figure}[p]
\centering
\begin{tikzpicture}
\begin{groupplot}[
 group style={group size=1 by 3,vertical sep=0.9cm},
 width=.88\textwidth,height=4.5cm,
 xmin=450,xmax=550,ymin=0,ymax=1.12,
 xtick={475,500,525},ytick={0,.25,.5,.75,1},grid=both,
 tick label style={font=\scriptsize},label style={font=\small},title style={font=\small},
 ylabel={Pico (V)},xlabel={Frequência (Hz)}]
\nextgroupplot[title={$\mu=0{,}5$: bandas de 0,25 V}]
\addplot[azul,very thick,ycomb,mark=*] table[x=frequencia_hz,y=amplitude_v,col sep=comma]{referencia_analitica/espectro_teorico_mu05.csv};
\nextgroupplot[title={$\mu=1{,}0$: bandas de 0,50 V}]
\addplot[azul,very thick,ycomb,mark=*] table[x=frequencia_hz,y=amplitude_v,col sep=comma]{referencia_analitica/espectro_teorico_mu10.csv};
\nextgroupplot[title={$\mu=1{,}5$: bandas de 0,75 V}]
\addplot[azul,very thick,ycomb,mark=*] table[x=frequencia_hz,y=amplitude_v,col sep=comma]{referencia_analitica/espectro_teorico_mu15.csv};
\end{groupplot}
\end{tikzpicture}
\caption{Espectros unilaterais \textbf{teóricos} da Equação~\eqref{eq:spectrum}. A linha central não muda; as laterais crescem proporcionalmente a $\mu$. Não são FFTs do TyphoonSim.}
\label{fig:espectro_teoria}
\end{figure}

\clearpage
\section{Resultados obtidos no TyphoonSim}
\ifcapturavhil
\input{resultados_vhil/medidas.tex}
O arquivo de dados contém \CapturaN\ amostras, capturadas em \texttt{\CapturaData}, a aproximadamente \CapturaTaxa\ amostras por segundo. O SHA-256 do modelo é \texttt{\CapturaHashModelo}; o do CSV bruto é \texttt{\CapturaHashDados}. A Tabela acima foi calculada automaticamente do CSV capturado. As frequências e amplitudes podem ser comparadas diretamente com a Tabela~\ref{tab:teoria}; o RMS do resíduo testa se as três componentes bastam para reconstruir o sinal do modelo ideal.

\begin{figure}[htbp]
\centering
\begin{tikzpicture}
\begin{groupplot}[
 group style={group size=1 by 3,vertical sep=.75cm},
 width=.85\textwidth,height=4.3cm,
 xmin=0,xmax=40,ymin=-2.7,ymax=2.7,
 xtick={0,10,20,30,40},grid=both,
 tick label style={font=\scriptsize},title style={font=\small},ylabel={V},xlabel={Tempo (ms)}]
\nextgroupplot[title={$\mu=0{,}5$}]
\addplot[azul,thin] table[x=tempo_ms,y=sinal_am,col sep=comma]{resultados_vhil/tempo_mu05.csv};
\addplot[laranja,thick] table[x=tempo_ms,y=env_pos,col sep=comma]{resultados_vhil/tempo_mu05.csv};
\addplot[laranja,thick] table[x=tempo_ms,y=env_neg,col sep=comma]{resultados_vhil/tempo_mu05.csv};
\nextgroupplot[title={$\mu=1{,}0$}]
\addplot[azul,thin] table[x=tempo_ms,y=sinal_am,col sep=comma]{resultados_vhil/tempo_mu10.csv};
\addplot[laranja,thick] table[x=tempo_ms,y=env_pos,col sep=comma]{resultados_vhil/tempo_mu10.csv};
\addplot[laranja,thick] table[x=tempo_ms,y=env_neg,col sep=comma]{resultados_vhil/tempo_mu10.csv};
\nextgroupplot[title={$\mu=1{,}5$}]
\addplot[azul,thin] table[x=tempo_ms,y=sinal_am,col sep=comma]{resultados_vhil/tempo_mu15.csv};
\addplot[laranja,thick] table[x=tempo_ms,y=env_pos,col sep=comma]{resultados_vhil/tempo_mu15.csv};
\addplot[laranja,thick] table[x=tempo_ms,y=env_neg,col sep=comma]{resultados_vhil/tempo_mu15.csv};
\addplot[vermelho,densely dashed,thick] table[x=tempo_ms,y=env_abs,col sep=comma]{resultados_vhil/tempo_mu15.csv};
\end{groupplot}
\end{tikzpicture}
\caption{Formas de onda \textbf{capturadas no TyphoonSim}: sinal AM, envoltória assinada e módulo da envoltória no regime sobremodulado.}
\label{fig:tempo_vhil}
\end{figure}

\begin{figure}[htbp]
\centering
\begin{tikzpicture}
\begin{groupplot}[
 group style={group size=1 by 3,vertical sep=.75cm},
 width=.85\textwidth,height=4.0cm,
 xmin=440,xmax=560,ymin=0,ymax=1.13,
 xtick={475,500,525},ytick={0,.25,.5,.75,1},grid=both,
 tick label style={font=\scriptsize},title style={font=\small},ylabel={Pico (V)},xlabel={Frequência (Hz)}]
\nextgroupplot[title={$\mu=0{,}5$}]
\addplot[azul,thick] table[x=frequencia_hz,y=amplitude_v,col sep=comma]{resultados_vhil/espectro_mu05.csv};
\nextgroupplot[title={$\mu=1{,}0$}]
\addplot[azul,thick] table[x=frequencia_hz,y=amplitude_v,col sep=comma]{resultados_vhil/espectro_mu10.csv};
\nextgroupplot[title={$\mu=1{,}5$}]
\addplot[azul,thick] table[x=frequencia_hz,y=amplitude_v,col sep=comma]{resultados_vhil/espectro_mu15.csv};
\end{groupplot}
\end{tikzpicture}
\caption{Espectros \textbf{obtidos por FFT das capturas VHIL}, com janela de Hann. O alargamento local dos picos é efeito da janela; a tabela de amplitudes usa o ajuste harmônico.}
\label{fig:espectro_vhil}
\end{figure}
\else
\noindent\fcolorbox{vermelho}{red!4}{\parbox{.93\textwidth}{\textbf{Captura pendente.} O Typhoon HIL não estava instalado no ambiente de autoria. Execute \arquivo{capturar_vhil.py} e depois \arquivo{analisar_captura.py} no computador com TyphoonSim. Os programas criarão \arquivo{resultados_vhil/medidas.tex} e os CSVs de figura; ao recompilar este mesmo fonte no Overleaf, esta seção trocará automaticamente a nota por tabela, formas de onda e FFTs reais. Esta prévia não deve ser apresentada como relatório experimental concluído.}}
\par\medskip
Enquanto essa execução não é feita, as Figuras~\ref{fig:tempo_teoria} e~\ref{fig:espectro_teoria} servem apenas como valores de referência para a comparação posterior. O professor solicitou explicitamente resultados obtidos na demonstração simulada, portanto essa etapa é indispensável para fechar a entrega.
\fi

\section{Discussão}
As previsões analíticas mostram duas mudanças distintas ao elevar $\mu$. No domínio da frequência, a amplitude de cada banda lateral passa de \SI{0.25}{\volt} para \SI{0.5}{\volt} e \SI{0.75}{\volt}, enquanto a portadora permanece com \SI{1}{\volt}. No tempo, o mínimo de $u$ passa de $0{,}5$ a $0$ e então a $-0{,}5$. A transição de sinal da envoltória assinada no terceiro regime explica por que uma detecção por envoltória baseada em $|u|$ distorce a mensagem. O espectro do sinal AM ideal, porém, continua restrito às três frequências da Equação~\eqref{eq:spectrum}; não se deve confundir a distorção do detector com criação de novas bandas no modulador linear.

\ifcapturavhil
Na captura real, a tabela de resultados permite testar as três relações: $A_{475}\approx A_{525}$, $A_{500}\approx\SI{1}{\volt}$ e $\widehat\mu\approx\mu$. A concordância entre as linhas previstas e os picos das FFTs associa a saída simulada à expansão trigonométrica; o cruzamento de zero de $u$ na última forma de onda associa a sobremodulação à inversão de fase. Pequenas diferenças numéricas podem decorrer do passo discreto, da janela e da representação em ponto flutuante, mas os limites de coerência foram verificados pelo script antes de gerar a tabela.
\else
A confirmação experimental dessas relações permanece pendente. Não foram atribuídos valores simulados às previsões da Tabela~\ref{tab:teoria} nem às figuras analíticas.
\fi

\section{Reprodutibilidade e entrega}
O pacote contém o \arquivo{.tse}, os scripts de captura e análise, este fonte \LaTeX{} e os CSVs \emph{teóricos} usados apenas na exposição. Após a aquisição, devem ser adicionados o CSV bruto VHIL, o JSON de metadados e a pasta \arquivo{resultados_vhil}. Os dados de Capítulo 3 não são reutilizados: a orientação exige implementar um novo modulador para a Equipe 04. O arquivo PDF de submissão final deve chamar-se \arquivo{Equipe_04_ModulacaoAM.pdf}, acompanhado dos scripts e dados auxiliares novos, conforme a orientação do professor.

\section{Conclusão}
Derivou-se a AM com portadora presente e construiu-se um esquema TyphoonSim que percorre três índices de modulação. A teoria prevê bandas laterais em \SI{475}{\hertz} e \SI{525}{\hertz}, de amplitudes proporcionais a $\mu$, e esclarece que a sobremodulação se identifica pelo sinal negativo da envoltória assinada, não por novas frequências no modulador ideal.
\ifcapturavhil
As medições VHIL incorporadas nas seções anteriores verificam essa previsão nos domínios do tempo e da frequência, com CSV bruto, taxa de amostragem e hashes registrados para reprodução.
\else
A etapa experimental permanece por executar no computador com Typhoon HIL. Até que os dados VHIL sejam incorporados, esta é uma prévia técnica e não a versão final para submissão.
\fi

\section*{Referências}
\begin{enumerate}[leftmargin=2.0em]
\item ECOM060 -- INSTRUMENTAÇÃO ELETRÔNICA. \emph{Orientação -- Capítulo 5: Análise Espectral, Sinais Modulados e Aleatórios}. Universidade Federal de Alagoas, semestre 2026.2, 3 p.
\item AGUIRRE, Luís Antonio. \emph{Fundamentos de Instrumentação}. São Paulo: Pearson Education do Brasil, 2013. Capítulo 5, seção 5.2. ISBN 978-85-8143-183-3.
\item TYPHOON HIL. \emph{HIL API: funções de carga, captura e leitura de sinais}. Documentação oficial. Disponível em: \href{https://www.typhoon-hil.com/documentation/typhoon-hil-api-documentation/hil_api.html}{API HIL do Typhoon}. Acesso em: 18 set. 2026.
\item TYPHOON HIL. \emph{Schematic Editor API: compilação de modelos}. Documentação oficial. Disponível em: \href{https://www.typhoon-hil.com/documentation/typhoon-hil-api-documentation/schematic_api.html}{API do Schematic Editor}. Acesso em: 18 set. 2026.
\end{enumerate}

\end{document}
